\documentclass[11pt,a4paper]{article}
\usepackage[utf8]{inputenc}
\usepackage[T1]{fontenc}
\usepackage[brazilian]{babel}
\usepackage{geometry}
\usepackage{longtable}
\usepackage{booktabs}
\usepackage{array}
\usepackage{titling}

\geometry{margin=2.2cm}

\title{Verificação e Validação de Software \\ \large Trabalho 2: Cobertura de Código e Teste de Mutação da Tarifa de Estacionamento}
\author{Eduardo Alcaria \\ Arthur Pimentel \\ Ethan Soares}
\date{}

\begin{document}
\maketitle
\vspace{-2.5em}

\section*{Introdução}

Este relatório aplica cobertura de código (JaCoCo 0.8.12) e teste de mutação
(PIT 1.17.4 com \texttt{pitest-junit5-plugin} 1.2.1) sobre a classe
\texttt{CalculadoraTarifa} e a suíte JUnit 5 produzida no Trabalho 1. O ambiente
usado foi Java 17 e Maven 3.9. A suíte original possui 25 execuções de teste
(\texttt{CalculadoraTarifaTest}); os casos novos deste trabalho estão na classe
driver \texttt{CalculadoraTarifaComplementarTest}, com 27 execuções, totalizando
52 execuções, todas aprovadas.

\section{Análise crítica da cobertura de código (JaCoCo)}

\subsection{Resultado inicial}

\begin{center}
\begin{tabular}{lrrr}
\toprule
Métrica & Perdidos & Cobertos & Cobertura \\
\midrule
Instruções & 0 & 160 & 100\% \\
Branches & 0 & 20 & 100\% \\
Linhas & 0 & 38 & 100\% \\
Métodos & 0 & 9 & 100\% \\
Complexidade ciclomática & 0 & 19 & 100\% \\
\bottomrule
\end{tabular}
\end{center}

Cobertura por método (todos com 0 itens perdidos): \texttt{calcular} (31
instruções, 2 branches), \texttt{calcularPerNoite} (16, 2),
\texttt{calcularDescontoVip} (8, 2), \texttt{calcularPorPermanencia} (31, 4),
\texttt{contarPernoites} (24, 2), \texttt{validarEntrada} (11, 2),
\texttt{validarSaida} (23, 6), além do construtor e do inicializador estático.

\subsection{Análise crítica}

A suíte do Trabalho 1 já atinge 100\% em todas as métricas do JaCoCo. Portanto
\textbf{não há linha, branch ou método sem cobertura} e nenhum caso de teste é
estritamente necessário para elevar a cobertura. Isso decorre de a suíte ter sido
projetada por particionamento em classes de equivalência e análise de valor
limite sobre cada regra de negócio, o que exercita naturalmente os dois lados de
cada decisão.

Cobertura de 100\% não prova que os testes verifiquem o resultado correto, apenas
que o código foi executado. Por isso foram acrescentados casos complementares,
orientados à técnica de projeto (e não a uma lacuna do JaCoCo), em regiões onde a
suíte original era rasa:

\begin{itemize}
\item Cliente VIP exercitado apenas em 3 casos; faltavam fronteiras de 60, 61 e 121 minutos e múltiplas diárias.
\item Pernoite testado com 1 e 2 diárias; faltava 3 diárias, virada de mês, virada de ano e ano bissexto.
\item Janela de saída proibida testada só nos limites 02:00 e 07:59 e no ponto 05:00; faltavam 02:01 e 07:58.
\item Permanência longa no mesmo dia testada até 4 h; faltavam 180/181 min e o teto de 23:59.
\item Saída às 01:59 de dois dias depois (uma pernoite, não duas).
\end{itemize}

\subsection{Novos casos de teste (resultados válidos)}

\small
\begin{longtable}{@{}l p{4.2cm} p{3.0cm} c r p{2.4cm}@{}}
\toprule
ID & Entrada & Saída & VIP & Esperado & Técnica \\
\midrule
\endhead
CC01 & 10/03/2026 08:00 & 10/03/2026 11:00 & não & 25,00 & Valor limite \\
CC02 & 10/03/2026 08:00 & 10/03/2026 11:01 & não & 30,00 & Valor limite \\
CC03 & 10/03/2026 08:00 & 10/03/2026 23:59 & não & 90,00 & Valor limite \\
CC04 & 10/03/2026 22:00 & 10/03/2026 23:59 & não & 20,00 & Valor limite \\
CC05 & 10/03/2026 23:00 & 11/03/2026 00:00 & não & 15,00 & Valor limite \\
CC06 & 10/03/2026 23:59 & 11/03/2026 01:59 & não & 20,00 & Valor limite \\
CC07 & 10/03/2026 23:59 & 11/03/2026 08:00 & não & 50,00 & Valor limite \\
CC08 & 10/03/2026 20:00 & 11/03/2026 08:01 & não & 50,00 & Valor limite \\
CC09 & 10/03/2026 20:00 & 12/03/2026 01:59 & não & 50,00 & Classe equiv. \\
CC10 & 10/03/2026 10:00 & 13/03/2026 10:00 & não & 150,00 & Classe equiv. \\
CC11 & 31/03/2026 20:00 & 01/04/2026 08:00 & não & 50,00 & Valor limite \\
CC12 & 31/12/2026 20:00 & 01/01/2027 08:00 & não & 50,00 & Valor limite \\
CC13 & 28/02/2028 20:00 & 01/03/2028 08:00 & não & 100,00 & Valor limite \\
CC14 & 10/03/2026 08:00 & 10/03/2026 09:00 & sim & 7,50 & Valor limite \\
CC15 & 10/03/2026 08:00 & 10/03/2026 09:01 & sim & 10,00 & Valor limite \\
CC16 & 10/03/2026 08:00 & 10/03/2026 10:01 & sim & 12,50 & Valor limite \\
CC17 & 10/03/2026 08:00 & 10/03/2026 08:20 & sim & 0,00 & Valor limite \\
CC18 & 10/03/2026 20:00 & 12/03/2026 08:00 & sim & 50,00 & Classe equiv. \\
CC19 & 10/03/2026 10:00 & 13/03/2026 10:00 & sim & 75,00 & Classe equiv. \\
\bottomrule
\end{longtable}
\normalsize

\subsection{Novos casos de teste (resultados inválidos, esperada \texttt{IllegalArgumentException})}

\small
\begin{longtable}{@{}l p{4.2cm} p{4.2cm} c p{2.8cm}@{}}
\toprule
ID & Entrada & Saída & VIP & Motivo \\
\midrule
\endhead
CI01 & 10/03/2026 20:00 & 11/03/2026 02:01 & não & Saída na janela proibida \\
CI02 & 10/03/2026 20:00 & 11/03/2026 07:58 & não & Saída na janela proibida \\
CI03 & 10/03/2026 20:00 & 12/03/2026 05:00 & não & Saída na janela proibida \\
CI04 & 10/03/2026 07:59 & 10/03/2026 09:00 & sim & Entrada antes das 08:00 \\
CI05 & 10/03/2026 20:00 & 11/03/2026 03:00 & sim & Saída na janela proibida \\
CI06 & 10/03/2026 10:00 & 10/03/2026 09:59 & sim & Saída antes da entrada \\
CI07 & 10/03/2026 10:00 & 10/03/2026 10:00 & sim & Saída igual à entrada \\
CI08 & 10/03/2026 10:00 & 09/03/2026 12:00 & não & Saída em dia anterior \\
\bottomrule
\end{longtable}
\normalsize

\subsection{Resultado após os novos casos}

Cobertura inalterada em 100\% (160/160 instruções, 20/20 branches, 38/38 linhas,
9/9 métodos), o que é esperado, já que não havia lacuna a fechar.

\section{Análise crítica do teste de mutação (PIT)}

\subsection{Resultado inicial}

Com o grupo de mutadores \texttt{DEFAULTS}:

\begin{center}
\begin{tabular}{lr}
\toprule
Mutantes gerados & 32 \\
Mortos & 30 \\
Sobreviventes & 2 \\
Mutation score & 93,75\% \\
Mutation score ajustado (sem equivalentes) & 100\% \\
\bottomrule
\end{tabular}
\end{center}

Para uma verificação mais rigorosa, a análise foi repetida com os grupos
\texttt{STRONGER} (42 mutantes: 39 mortos, 1 \emph{timeout}, 2 sobreviventes) e
\texttt{ALL} (94 mutantes: 88 mortos, 4 \emph{timeout}, 2 sobreviventes). Os
\emph{timeouts} ocorrem em \texttt{contarPernoites}: mutantes que tornam o
\texttt{while} infinito, contabilizados pelo PIT como detectados.

\subsection{Mutantes sobreviventes}

Os mesmos dois sobreviventes aparecem em todos os grupos, ambos em
\texttt{calcularPorPermanencia}, linha 53, \texttt{if (minutos <= MINUTOS\_TARIFA\_FIXA)}:

\begin{center}
\small
\begin{tabular}{@{}p{4.3cm} p{4.8cm} p{5.2cm}@{}}
\toprule
Mutador & Mutação & Classificação \\
\midrule
\texttt{ConditionalsBoundaryMutator} & \texttt{<=} vira \texttt{<} & Equivalente \\
\texttt{RemoveConditionalMutator\_ORDER\_ELSE} & condição substituída por \texttt{false} & Equivalente \\
\bottomrule
\end{tabular}
\end{center}

\textbf{Mutante 1 (\texttt{<=} para \texttt{<}).} A única entrada em que o
comportamento da condição muda é \texttt{minutos == 60}. No original retorna
\texttt{15,0}. No mutante, o fluxo segue: \texttt{minutosExcedentes = 0},
\texttt{horasAdicionais = (0 + 59) / 60 = 0}, resultado \texttt{15 + 0 * 5 = 15,0}.
Mesma saída, logo equivalente.

\textbf{Mutante 2 (condição sempre falsa).} Para \texttt{21 <= minutos <= 60}
(a cortesia já foi tratada na linha 50), \texttt{minutosExcedentes} está em
$[-39, 0]$, então \texttt{minutosExcedentes + 59} está em $[20, 59]$ e a divisão
inteira por 60 resulta em 0, gerando \texttt{15,0}, idêntico ao original. Logo
equivalente.

Como ambos são equivalentes, nenhum teste é capaz de matá-los sem alterar o
código de produção. Opcionalmente, o ramo da linha 53 poderia ser removido
(refatoração que preserva o comportamento e elimina a redundância), mas isso foge
ao escopo do trabalho, que consiste em complementar testes.

\subsection{Análise crítica}

A suíte do Trabalho 1 já mata todos os mutantes não equivalentes, inclusive com
mutadores além do padrão. Não há mutante não equivalente a cobrir com novos
casos. Os casos CC01 a CC19 e CI01 a CI08 reforçam as regiões de fronteira
descritas na Seção 1, reduzindo o risco de regressões não capturadas pelo
conjunto de mutadores do PIT (por exemplo, falhas de arredondamento em horas
adicionais e virada de calendário).

\subsection{Resultado após os novos casos}

\begin{center}
\begin{tabular}{lrr}
\toprule
Métrica & Antes & Depois \\
\midrule
Execuções de teste & 25 & 52 \\
Mutantes gerados (\texttt{DEFAULTS}) & 32 & 32 \\
Mortos (\texttt{DEFAULTS}) & 30 & 30 \\
Sobreviventes (\texttt{DEFAULTS}) & 2 & 2 \\
Mutation score (\texttt{DEFAULTS}) & 93,75\% & 93,75\% \\
Mutantes gerados (\texttt{ALL}) & 94 & 94 \\
Sobreviventes (\texttt{ALL}) & 2 & 2 \\
\bottomrule
\end{tabular}
\end{center}

O score não muda porque os únicos sobreviventes são equivalentes. Os novos casos
ampliam a verificação do comportamento sem alterar o resultado das métricas.

\section{Como reproduzir}

\begin{verbatim}
mvn test
mvn test-compile org.pitest:pitest-maven:mutationCoverage
mvn test-compile org.pitest:pitest-maven:mutationCoverage -Dmutators=ALL
\end{verbatim}

Relatório JaCoCo em \texttt{target/site/jacoco/index.html} e relatório PIT em
\texttt{target/pit-reports/index.html}.

\end{document}
